\subsection{B.4 supplement. A Hodge obstruction to the central strata}

\begin{theorem}[B.4.3; Known consequences with explicit derivation]
Let $g:\mathcal Z\to\Delta_s$ be a proper holomorphic map from a
smooth complex fourfold, smooth over $\Delta_s^*$, with reduced
normal-crossings central fibre $D$. Suppose its nearby fibres are
biholomorphic to the six-sphere members of the specified family,
after any finite base change $w=s^N$. Then
\[
 H^1(D;\C)=0,\qquad h^1(D,\OO_D)\geq1.
\]
Consequently the natural map
$H^1(D;\C)\to H^1(D,\OO_D)$ is not surjective. In particular,
$D$ cannot be the analytification of a proper complex algebraic
variety with normal-crossings singularities.

If $D$ has \emph{simple} normal crossings, at least one of its
nonempty closed strata fails the $\partial\bar\partial$-lemma.
In particular, no such simple normal-crossings filling can have
all irreducible components K\"ahler or all irreducible components
Moishezon. These obstructions hold for every $N\geq1$.
\end{theorem}
\begin{proof}
The fundamental-group surjection in B.4.2 gives
$\pi_1(D)=1$, and hence $H^1(D;\C)=0$. We check the other
input directly. The audited root-curve direct images in E.2 and
E.3, for $\mathfrak B=\PP(3,4)$, are
\[
 \bar f_*\OO_X=\OO_{\mathfrak B},\qquad
 R^1\bar f_*\OO_X=\OO_{\mathfrak B}\oplus
                          \OO_{\mathfrak B}(-1).
\]
The coarse curve is $\PP^1$, and taking invariants by the finite
stabilizers is exact in characteristic zero. Thus
$H^1(\mathfrak B,\OO)=H^2(\mathfrak B,\OO)=0$.
Moreover $H^0(\mathfrak B,\OO)=\C$ and
$H^0(\mathfrak B,\OO(-1))=0$: in the weighted homogeneous
coordinate ring the latter would require a monomial of weighted
degree $-1$. The degree-one Leray edge sequence therefore gives
\[
 H^1(X,\OO_X)\simeq
 H^0(\mathfrak B,R^1\bar f_*\OO_X)=\C.
\]
This confirms the precise R4 input, rather than inferring
$h^1(\OO)$ from tangent-bundle cohomology. The elementary weighted
degree counts are recorded in \path{code/b_hodge_obstruction.py}.
Upper semicontinuity for the proper flat map $g$ gives
$h^1(D,\OO_D)\geq1$; flatness also follows directly from the
local equation $s=z_1\cdots z_k$.

There is a family-independent check of the required lower bound.
For any compact complex threefold $X$ diffeomorphic to $S^6$,
$c_1(X)=c_2(X)=0$ in integral cohomology, so analytic
Hirzebruch--Riemann--Roch gives
$\chi(\OO_X)=\langle c_1c_2,[X]\rangle/24=0$.
The analytic index theorem used here applies without a K\"ahler
assumption; see \citet[Theorem 4.3]{AtiyahSinger1968IIIE}.
A holomorphic three-form is closed and, since $H^3(X;\C)=0$,
exact. Stokes' theorem applied to its product with its complex
conjugate, whose suitable constant multiple is pointwise
nonnegative, shows that it vanishes. Serre duality then gives
$h^3(\OO_X)=0$. It follows that
$h^1(\OO_X)=1+h^2(\OO_X)\geq1$.

For the algebraic assertion we use the precise properness
hypothesis in \citet[Section 1.D, p.120]{GrafKovacs2014Root}.
For a proper complex algebraic variety with Du Bois singularities,
the canonical map from singular cohomology to coherent cohomology
is surjective in every degree. Normal-crossings singularities are
Du Bois, including self-intersections of a global component.
Applied in degree one, this contradicts the two groups just
computed. No algebraicity of the total space is required for
this contradiction; algebraicity of the central fibre alone
would suffice.

For the simple normal-crossings assertion, suppose instead that
every nonempty intersection of irreducible components of $D$
is a $\partial\bar\partial$-manifold. This is exactly condition
(K1) in \citet[Definition 2.3]{Chen2026HodgeRoot}.
Theorem 2.4(4) of that source, which states the relevant
Steenbrink degeneration theorem under (K1), gives degeneration
at $E_1$ of the nearby Hodge--de Rham spectral sequence.
All its hypotheses have been specified: $g$ is proper,
$\mathcal Z$ is smooth, $D$ is reduced and simple normal
crossings, and its closed strata satisfy (K1). Neither a
K\"ahler metric on $\mathcal Z$ nor a local invariant-cycle
theorem is being assumed. Degeneration in total degree one
would imply
\[
 0=b_1(X_s)=h^0(X_s,\Omega^1_{X_s})+h^1(X_s,\OO_{X_s})
          \geq1,
\]
which is impossible.

Closed smooth submanifolds of a compact K\"ahler manifold are
K\"ahler. Closed strata of a Moishezon component are Moishezon:
a projective modification, followed by taking a component of
the inverse image that dominates the stratum and a general
relative linear section if needed, gives a generically finite
projective cover. Its meromorphic functions show that the
stratum has maximal algebraic dimension. Smooth Moishezon
manifolds satisfy the $\partial\bar\partial$-lemma, as recalled
in \citet[Section 2.1, p.8]{Chen2026HodgeRoot}. Thus either
of the two stated conditions on components would imply (K1).
Finally, finite base change does not change the biholomorphism
type or either cohomological value of a nonzero fibre, so none
of the contradictions depends on $N$.
\end{proof}

This is an application of existing Hodge-theoretic results, not
a claim of a new degeneration theorem. Searches on 15 November
2026 for six-sphere semistable degenerations, normal-crossings
limits, Moishezon central fibres, and Conjecture 8.2 with Hodge
theory did not locate this specific application. We opened
Chen's version 2 of 16 October 2026 and checked Definition 2.3
and Theorem 2.4(4), including the condition inserted before
parts (3)--(4); omitting that condition would make the
statement false. We also opened the original Graf--Kov\'acs
article and checked the proper Du Bois assertion in Section
1.D. The separate source audit records the searches.

The theorem prevents a repair made solely from compact
Moishezon components in a simple normal-crossings model.
It does not prove that an arbitrary analytic filling must have
such components. In particular, no argument here makes an
unknown three-dimensional component over the cusp Moishezon.
The existence and least base-change exponent in B.4.1 therefore
remain open in the unrestricted analytic category.
The Hopf-component example of
\citet[Theorem 1.4 and Example 2.1]{Nishiguchi1988IIRoot}
shows why an analytic degeneration cannot simply be assumed to
have class-$\mathcal C$ strata. The toric-Kato construction of
\citet[Theorem 7.1]{IstratiEtAl2022Root} is another relevant
non-K\"ahler analogue. Neither source supplies the required
cusp comparison for the six-sphere family.

\begin{proposition}[B.4.4 (Proved): extension of the marked fibration]
Suppose $g:\mathcal Z\to\Delta_s$ is a proper holomorphic map
from a smooth complex fourfold, with reduced normal-crossings
central fibre, and its restriction over $\Delta_s^*$ is
isomorphic to a finite base change of the specified six-sphere
family. The given marked holomorphic map
\[
 f^*:\mathcal Z|_{\Delta_s^*}\longrightarrow\PP^1_t
\]
extends to a meromorphic map $\mathcal Z\dashrightarrow\PP^1_t$.
Its graph gives a proper projective modification
$\mathcal Z'\to\mathcal Z$, an isomorphism over
$\Delta_s^*$, on which the marked map is holomorphic.
No assertion that $\mathcal Z'$ is smooth or semistable is
included in this statement.
\end{proposition}
\begin{proof}
We shrink the disc whenever necessary. Put
$L=K_{\mathcal Z/\Delta_s}^{-2}$. The relative canonical
bundle here is the line bundle
$K_{\mathcal Z}\otimes g^*K_{\Delta_s}^{-1}$; its
restriction to a nonzero fibre is that fibre's canonical bundle.
The precise R1 input is
$K_{X_s}^{-2}\simeq f_s^*\mathcal O_{\PP^1}(1)$,
with connected fibres and no further sections. Therefore
$h^0(X_s,L|_{X_s})=2$, and its complete linear system is
the given map $f_s$, up to a projective change of coordinate.
The three critical values are intrinsically distinguished by
their fibres: multiplicity three, multiplicity four, and reduced
normal crossings, respectively.

Grauert's proper direct-image theorem
\citep[Section 6, Hauptsatz I]{Grauert60push} makes
$\mathcal F=g_*L$ coherent. It is torsion-free: multiplication
by a nonzero base germ is injective on $L$, because the local
semistable equation makes $g$ flat, and hence is injective on
its direct image. A coherent torsion-free sheaf on the smooth
disc is locally free. Its rank is two. On the punctured disc,
constant $h^0=2$ and the analytic base-change theorem checked
in E.3 identify its fibres with the complete spaces of sections.
Choose a frame $(\sigma_0,\sigma_1)$ of $\mathcal F$.
Evaluation gives two global holomorphic sections of $L$ on
$\mathcal Z$, and thus a meromorphic map
$F_0=[\sigma_0:\sigma_1]$ to $\PP^1$.

For completeness its graph is analytic and proper without an
appeal to an unproved extension of a punctured graph. In a local
trivialization of $L$, write the sections as $a,b$. Remove their
common divisorial factor in the regular local ring of the smooth
space $\mathcal Z$. The equation
$v a-u b=0$ in $\mathcal Z\times\PP^1_{[u:v]}$
then has a unique component dominating $\mathcal Z$; it is
the closure of the graph where $(a,b)\ne(0,0)$. These local
components glue. Denote the resulting reduced graph by
$\mathcal G$. The projection $\mathcal G\to\mathcal Z$
is proper and projective, and is an isomorphism over
$\Delta_s^*$. The map
\[
 H:\mathcal G\longrightarrow\PP^1\times\Delta_s,
 \qquad H=(F_0,g),
\]
is proper. Locally $\mathcal G$ is a hypersurface in the smooth
space $\mathcal Z\times\PP^1$. The vanishing of the maximal
Jacobian minors for its equation and the two target coordinates
defines the closed analytic locus $C$ where $H$ is not a
submersion; singular points of $\mathcal G$ are included.

The image $H(C)$ is a closed analytic subset of the target.
One may deduce this precise proper-image assertion from the same
direct-image theorem: the coherent sheaf
$H_*\mathcal O_C$ has support exactly $H(C)$, since the
constant section $1$ is nonzero at every nonempty proper fibre.
Over $\Delta_s^*$ this image consists exactly of three distinct
critical-value sections. It cannot have a two-dimensional
component, since such an analytic component would contain an
open set of the connected two-dimensional target and hence meet
the punctured region in dimension two. Remove any component
contained in $\PP^1\times\{0\}$. The remaining reduced
analytic curve is proper with zero-dimensional fibres over the
disc, and is therefore finite over it.

Each of the three labelled sections on the punctured disc belongs
to a degree-one branch of that finite curve. Indeed their labels
are intrinsic and cannot be permuted by continuation. The
normalization of a degree-$d$ branch over a punctured disc has
cyclic monodromy of length $d$; a single-valued labelled section
therefore has $d=1$. The normalization extends it holomorphically
to a map from the full disc to $\PP^1$. Let these three maps be
$c_1(s),c_2(s),c_0(s)$, in the multiplicity-three,
multiplicity-four and cusp order. They may collide at $s=0$.

For $s\ne0$ the unique projective transformation sending
$(c_1,c_2,c_0)$ to $(0,1,\infty)$ has meromorphic coefficients
at $0$. In a common affine chart its value at $z$ is
\[
 A_s(z)=
 \frac{(z-c_1(s))(c_2(s)-c_0(s))}
      {(z-c_0(s))(c_2(s)-c_1(s))};
\]
homogeneous brackets give the same formula if a value is infinite.
The denominators are not identically zero, since the three values
are distinct on the punctured disc. Thus $A_s\circ F_0$ is
meromorphic on $\mathcal Z$. On every nonzero fibre it agrees
with the actual marked $f_s$, because both are the complete
bi-anticanonical map and have the same three labelled values.
This equality is an equality of maps over the punctured disc,
not just a fibrewise biholomorphism statement.

Finally the graph of $A_s\circ F_0$ is constructed by the same
local two-section argument after clearing powers of the base
coordinate in its meromorphic coefficients. Its projection to
$\mathcal Z$ is the asserted proper projective modification,
and its second projection is the required holomorphic marked map.
\end{proof}

The proposition removes a possible essential projective-coordinate
ambiguity. It does not extend a punctured isomorphism to any of
the auxiliary models of B.3, and does not make a component over
$t=\infty,s=0$ Moishezon. Those would be additional statements.

\begin{theorem}[B.4.5 (Proved): the obstruction on normalization components]
Let $D$ be a compact reduced complex space with normal-crossings
singularities, allowing self-intersections of global components.
Suppose that every connected component of its smooth normalization
satisfies the $\partial\bar\partial$-lemma. Then the natural map
\[
 H^1(D;\C)\longrightarrow H^1(D,\OO_D)
\]
is surjective.

Consequently every smooth-total-space reduced normal-crossings
filling of the specified six-sphere family, after any finite base
change, has a normalization component which fails the
$\partial\bar\partial$-lemma. In particular its normalization
cannot consist entirely of compact K\"ahler or Moishezon manifolds.
This conclusion includes normal crossings which are not simple.
\end{theorem}
\begin{proof}
We give a degree-one argument directly on $D$; no nearby
Hodge-de Rham degeneration theorem is needed. Let
$D^{[p]}$ be the smooth compact branch space whose points are
ordered $(p+1)$-tuples of distinct local branches through one
point of $D$. Equivalently, we normalize the loci where these
branches meet, retaining their ordering. These are distinct
\emph{local branches}, even when they belong to one global
irreducible component. The finite holomorphic maps
$\pi_p:D^{[p]}\to D$ and their face maps are specified by
forgetting one branch. The symmetric group acts by reordering;
we take the alternating summand, with the sign character, in
each degree. Averaging over this finite group is exact over
$\C$. Write $\delta$ for the alternating sum of the face
pullbacks.

Both augmented complexes
\[
 \C_D\longrightarrow
 \bigl((\pi_p)_*\C_{D^{[p]}}\bigr)_{p\geq0}^{\mathrm{alt}},
 \qquad
 \OO_D\longrightarrow
 \bigl((\pi_p)_*\OO_{D^{[p]}}\bigr)_{p\geq0}^{\mathrm{alt}}
\]
are resolutions. This can be checked locally without a global
ordering. In a chart $z_1\cdots z_k=0$, the constant complex is
the augmented cochain complex of the simplex on the $k$ branches.
For the holomorphic complex, fix a monomial in the convergent
power series ring. It survives on exactly the intersections
indexed by subsets of those $i$ for which its $z_i$-exponent is
zero. If there are no such indices it already vanishes on $D$;
otherwise its coefficient complex is again an augmented simplex
complex. The usual contraction is a finite combination of
coordinate restrictions, so preserves convergence. This proves
exactness for analytic functions, not just for formal series.
The descriptions agree under permutations of the local branches.
The finite simplex and permutation identities for the at most
four branches in our threefold application are checked exactly
in \path{code/b_nc_degree_one.py}.

Resolve the holomorphic terms by smooth Dolbeault forms and
the constant terms by smooth de Rham forms on the branch spaces.
These compute the respective hypercohomologies. Indeed finite
pushforward is exact here, and partitions of unity from local
ambient smooth functions pull back to make the resulting sheaves
fine. The ordinary de Rham lemma and Dolbeault lemma apply on
each smooth branch space. Projection to the $(0,\bullet)$
part defines the natural comparison between the two total
complexes and induces $H^1(D;\C)\to H^1(D,\OO_D)$.

Represent a coherent degree-one class by a total cocycle
$(\alpha,f)$, where $\alpha$ is a smooth $(0,1)$-form on
$D^{[0]}$ and $f$ is an alternating smooth function on
$D^{[1]}$. Choose the total-complex signs so that its equations
are
\[
 \bar\partial\alpha=0,\qquad
 \delta\alpha=\bar\partial f,\qquad \delta f=0.
\]
On a normalization component, $\partial\alpha$ is of type
$(1,1)$, is $\partial$-exact and is $\bar\partial$-closed.
The $\partial\bar\partial$-lemma gives a smooth function $h$
with $\partial\alpha=\partial\bar\partial h$.
Do this on every component, and put
\[
 \alpha'=\alpha-\bar\partial h,\qquad
 f'=f-\delta h.
\]
The pair represents the same coherent class and still satisfies
the displayed cocycle equations; now $d\alpha'=0$.
Holomorphic face maps commute with both differential operators,
so
\[
 \partial\bar\partial f'
       =\partial\delta\alpha'=\delta\partial\alpha'=0.
\]
A smooth complex-valued pluriharmonic function on a compact
connected complex manifold is constant: apply the maximum
principle to its real and imaginary parts in local complex
coordinates. Therefore $f'$ is constant on each connected
component of $D^{[1]}$. It follows that
$df'=0$ and $\delta\alpha'=0$.

Thus the same pair $(\alpha',f')$, with $\alpha'$ viewed as
a complex-valued smooth one-form, is a cocycle of the total
de Rham branch complex. Its image is the original coherent
class. This proves surjectivity. Notice that no
$\partial\bar\partial$ hypothesis on $D^{[1]}$ or higher
branch spaces was used; compactness suffices for the only
function argument there.

For the asserted application, B.4.2 and B.4.3 give
$H^1(D;\C)=0$ and $h^1(D,\OO_D)\geq1$.
Surjectivity is impossible, so some normalization component
fails the lemma. Compact K\"ahler and smooth Moishezon
manifolds satisfy it, as in the primary comparison checked
in B.4.3. The same argument applies after every finite base
change, since its smooth fibres are still the specified
six-spheres.
\end{proof}

This strengthens the simple-normal-crossings obstruction in
B.4.3 by using only degree one. It also avoids an unverified
extension of an algebraic Du Bois theorem to arbitrary analytic
spaces. Searches for normal-crossings cohomology with K\"ahler
normalization, the normalization Dolbeault complex, and
degree-one surjectivity located the broader Hodge-theoretic
setting discussed by \citet{Chen2026HodgeRoot} and the proper
algebraic statement of \citet{GrafKovacs2014Root}; we do not
claim priority for this elementary consequence. The proof above
states and verifies the particular analytic comparison required
here. It still does not prove that all possible analytic
normalization components satisfy its hypothesis.

\begin{proposition}[B.4.6; Proved; a restricted non-K\"ahler cap obstruction]\label{push:B46}
Let $D$ be a compact reduced normal-crossings complex
space. Suppose its irreducible components have been
partitioned into closed unions $D_0,\ldots,D_s$, with
their reduced induced structures. Assume that
\[
 H^1(D_i,\mathbb C)\longrightarrow H^1(D_i,\mathcal O)
 \quad\hbox{is surjective for every }i,
 \tag{B.46a}
\]
and that the analogous map is injective for every
pairwise intersection $D_i\cap D_j$. Then
\[
 H^1(D,\mathbb C)\longrightarrow H^1(D,\mathcal O_D)
 \quad\hbox{is surjective}.
 \tag{B.46b}
\]

In particular these hypotheses hold for the following
specified cap assembly. The horizontal block $D_0$ is
a two-section pinching of a ruled bundle over an
analytic locally trivial torsor $V\to C$ of B.3.8,
with conductor map $A$ covering $\sigma$ and
$A^{12}=t_{6P}$. The normalization may undergo the
modifications of B.3.7, or any smooth bimeromorphic
modifications for which the two conductor restriction
maps on $H^1$ are the same pullbacks as before.
The other blocks are smooth hyperbolic toric Kato
threefolds or compact smooth spaces satisfying
(B.46a), for example class-$\mathcal C$ spaces.
Assume every pairwise overlap is either a reduced
union of invariant divisors of a hyperbolic toric
Kato threefold, or a compact normal-crossings surface
whose normalized components and normalized double
curves are rational. No such assembly can be a
semistable central fibre of the $S^6$ family.

This excludes the stated cap construction even when
an invariant Kato surface in an overlap kills a
fundamental-group generator of a Kato component.
It is not a classification of arbitrary analytic
cusp components or of their possible overlaps.
\end{proposition}
\begin{proof}
Write $D^{[p]}$ for the disjoint union of intersections
of $p+1$ distinct blocks. The augmented closed-cover
complexes resolve both $\mathbb C_D$ and $\mathcal O_D$.
For the coherent assertion this can be checked in a
normal-crossings coordinate chart. The local equation
of $D$ is a product of distinct coordinates, and each
block groups a disjoint subset of those factors.
The corresponding ideals are generated by products
in disjoint sets of variables. The ordinary alternating
intersection complex is exact, as one sees on each
monomial, where it is the augmented complex of a
simplex. The same contraction acts coefficientwise
on convergent power series. For constant sheaves,
small intersections are contractible unions of
coordinate subspaces, and the same stalkwise simplex
complex is exact. This proves both resolutions also
when the global blocks are self-crossing.

Their hypercohomology spectral sequences have
\[
 E_1^{p,q}(\mathcal F)=H^q(D^{[p]},\mathcal F),
 \qquad \mathcal F=\mathbb C\ \hbox{or}\ \mathcal O.
\]
Every connected compact reduced branch intersection
has only constant holomorphic functions. Thus the
comparison maps in the entire $q=0$ row are
isomorphisms. They induce isomorphisms on
$E_2^{1,0}$ and $E_2^{2,0}$. Surjectivity in (B.46a)
and injectivity on pairwise intersections imply
surjectivity on $E_2^{0,1}$: lift a coherent
degree-one class on the blocks, then its topological
restriction difference vanishes because its coherent
image vanishes and the pairwise comparison is
injective.

Now use the natural five-term exact sequences
\[
 0\longrightarrow E_2^{1,0}
 \longrightarrow H^1(D,\mathcal F)
 \longrightarrow E_2^{0,1}
 \xrightarrow{d_2} E_2^{2,0}.
\]
The isomorphism on the last term makes a lift of a
coherent $d_2$-closed element topologically
$d_2$-closed. The isomorphism on the first term then
gives (B.46b). If the block comparisons are
isomorphisms, the same argument also proves
injectivity, provided the map on $E_2^{0,1}$ is
injective.

We verify the horizontal-block assertion rather than
apply the pairwise criterion directly to its two
conductor branches. By B.3.8,
$H^1(V,\mathbb C)=H^1(C,\mathbb C)$ and
$H^1(V,\mathcal O)=H^1(C,\mathcal O_C)$.
The first equality is the Leray calculation there;
the second also follows from
$R^1\pi_*\mathcal O_V=\mathcal O_C(-\infty)$,
which has no global section. The action of $A$ on
these groups is $\sigma^*$, since it covers $\sigma$
and the displayed maps are pullbacks from the base.
The characteristic polynomial
\[
 p(x)=\Phi_{12}(x)\Phi_6(x),\qquad p(1)=1,
\]
shows that $1-A^*$ is invertible on both groups.
The normalization sequence of the two-section
pinching, for constants and for holomorphic
functions, therefore gives
\[
 H^1(D_0,\mathbb C)=\mathbb C,
 \qquad H^1(D_0,\mathcal O)=\mathbb C.
\]
Both classes come from the same difference map on
constant functions of the conductor. Their natural
comparison is an isomorphism. Projective-bundle
pullback identifies the normalization groups with
those of $V$. The stipulated modifications preserve
these groups and both restriction maps; in particular
the explicit modifications of B.3.7 have this property.
This grouped computation is essential: on an individual
conductor $V$, the comparison has dimensions six and
three and is not injective.

For a hyperbolic toric Kato manifold $Z$, its
fundamental group is $\mathbb Z$, as in
\cite[Section 2, p.1352]{IstratiEtAl2022Root};
\cite[Proposition 8.4 and Theorem 8.5]{IstratiEtAl2022Root}
give $H^0(Z,\Omega_Z^1)=0$ and
$H^1(Z,\mathcal O_Z)=\mathbb C$.
The exact sequence
$0\to\mathbb C\to\mathcal O\xrightarrow{d}d\mathcal O\to0$
shows that $H^1(Z,\mathbb C)\to H^1(Z,\mathcal O)$
is injective, because a global section of
$d\mathcal O$ is a closed holomorphic one-form.
The two dimensions are one, so it is an isomorphism.
The same reasoning applies to hyperbolic toric Kato
surfaces. Compact smooth rational varieties have
both groups zero. Class-$\mathcal C$ blocks satisfy
(B.46a) by their $\partial\bar\partial$ property.

Finally let $F$ be a reduced union of invariant
divisors in a hyperbolic toric Kato threefold.
Its normalized components are compact toric surfaces
or hyperbolic toric Kato surfaces
\cite[proof of Theorem 8.8, p.1380]{IstratiEtAl2022Root}.
Its normalized double curves are rational, since
every invariant curve is rational in the hyperbolic
case \cite[Theorem 5.2]{IstratiEtAl2022Root}.
Use the ordered-local-branch resolution from B.4.5.
The comparisons on component $H^1$ are isomorphisms,
and the groups on normalized double curves vanish.
The same five-term argument shows that
$H^1(F,\mathbb C)\to H^1(F,\mathcal O_F)$ is an
isomorphism. For an overlap with rational normalized
components and double curves, the argument is
identical with zero component $H^1$. Thus every
pairwise overlap in the statement has the required
injective comparison, proving (B.46b) for the
specified assembly.

For an actual semistable $S^6$ centre, B.4.2 gives
$H^1(D,\mathbb C)=0$, while upper semicontinuity
and the nearby value $h^1(\mathcal O)=1$ in B.4.3
give $h^1(D,\mathcal O_D)\geq1$. This contradicts
(B.46b). The obstruction uses the cohomology maps
on the overlaps; it does not assume that every
Kato component's loop survives in the union.
\end{proof}


\paragraph{The remaining analytic comparison.}
We also examined whether general results on degenerations of tori
would force an arbitrary filling into the excluded toroidal class.
The distinction between a construction theorem and such a comparison
theorem is essential. \citet[Theorem 5.1.3 and Sections 5.2--5.4]
{KKN2008Root} construct proper analytic models of log complex tori
without assuming a polarization. Their input is a specified log
complex torus and an admissible pairing; the theorem does not assert
that every proper analytic filling is one of these models. We have
not constructed that global log input with the required finite-place
and cusp markings for the six-sphere family.

The classification in \citet[Theorem 3.1]{KamenovaToriRoot} assumes
a K\"ahler total space and trivial canonical bundle. Those hypotheses
have not been proved for a slice of an arbitrary filling here.
\citet[Section 3, Proposition 3.1]{SchroeerSiebert2002Root} explicitly
study normal-crossings degenerations with Hopf normalization, but
their smooth fibres are primary Kodaira surfaces. This is a source
of alternative analytic models to test, not a comparison with our
punctured family.

Even extending the marked map in B.4.4 does not extend a chosen
punctured isomorphism to a second model. Properness alone cannot
justify that inference: the automorphism
$(s,[z_0:z_1])\mapsto(s,[e^{1/s}z_0:z_1])$ of
$\Delta^*\times\PP^1$ has no meromorphic extension across zero,
although both product families are proper over the disc. Its
restriction to $[1:1]$ has an essential singularity. Thus an
unrestricted proof still needs a specific extension or boundedness
argument for the comparison, or an actual non-toroidal cusp
construction with the prescribed punctured restriction.

